\section{Limitations}
\label{sec:limitations}

\paragraph{Coverage and selection.}
The \StudyBenchmarkCount{} reported benchmarks were selected from 554 candidates through evidence,
scope, testability, and evaluation-frame review, not sampled randomly. The candidate labels overlap,
and the collection includes coding benchmarks outside the study scope. Benchmarks with accessible data,
executable environments, and automated scoring are easier to include; hardware-dependent,
embodied, and human-judged tasks are underrepresented. This selection may favor tasks that a
coding agent can address through files and programs. Domain-level counts describe this suite
and do not estimate performance over all tasks in a domain.

\paragraph{One agent and incomplete controlled comparisons.}
The direct-Codex arm uses a particular client, backbone, reasoning effort, and pass@1 protocol
(Appendix~\ref{sec:appendix:setup}). Results do not establish the best achievable performance
of either coding agents or specialized systems. The 25\%-subset comparator campaign is still
running; preliminary runs, full-set historical measurements, and incomplete subsets cannot be
combined as if they shared a denominator. Reference-informed and tool-library ablations remain
pending. We therefore do not yet claim a general causal explanation for the observed gaps.

\paragraph{Published scores and statistical interpretation.}
A paper score can differ in backbone, dataset version, judge, resources, and inference budget.
It is an external reference, not a shared-backbone replication. A confidence interval around
Codex alone does not account for uncertainty in the published score, and failure to detect a
difference does not establish equivalence. The benchmarks use different metrics, so we report
per-benchmark results and explicit counting denominators rather than an average of raw scores.
Multiple comparisons and single-run stochasticity further limit inferential claims.

\paragraph{Protocol repairs and resource access.}
Several initial Codex evaluations did not supply the intended environment or working web
access. Updated runs replace earlier results where available; six scored variants remain
explicitly labelled, while RareBench still lacks a valid network-enabled measurement. ChemBench's
150-question open-web pilot is an intentionally scoped result, not a completed full-benchmark
run. SSTQA's available score is an English-split judge variant; its Chinese-split primary
evaluation remains pending. Comparator patches, unavailable case banks,
different judges, and answer-key revisions are documented per row. Giving both arms the
benchmark's environment does not eliminate differences in a comparator's own retrieval or tool
implementation; those components are part of the method under study.

\paragraph{Audit validity and exposure.}
AI-assisted extraction and adversarial review can both make mistakes. The 2026-09-11 audit
targeted 25 previously unassessed bindings and excluded five already verified bindings, so its
finding rate is not a prevalence estimate for the 30 benchmarks in the task inventory or the literature.
Findings describe artifacts at a particular date; their severity and expected score direction
are not measured treatment effects. SOTA-finder accuracy against an independently adjudicated
reference set remains pending. Isolation and trajectory scans reduce runtime answer leakage
but cannot exclude pretraining contamination.
